\documentclass[12pt]{article}
\usepackage[utf8]{inputenc}
\usepackage{geometry}
\usepackage{hyperref}
\usepackage{enumitem}
\usepackage{setspace}
\geometry{margin=1in}
\setlist{topsep=4pt,itemsep=2pt,leftmargin=*}
\setlength{\parskip}{6pt}
\setlength{\parindent}{0pt}

\title{Ultimate Guide to Vibe Coding V1.2.2}
\author{Nicolas Zullo \\ \href{https://x.com/NicolasZu}{https://x.com/NicolasZu}}
\date{Creation Date: March 12, 2025 \\ Last Update Date: January 15, 2026}

\begin{document}

\maketitle
\bigskip
\hrule
\bigskip

\section*{Key Principle}
\textbf{Planning is everything.} Do not let the AI plan autonomously, or your codebase will become an unmanageable mess.

\section{Getting Started}
To begin vibe coding, you only need one of these two tools:
\begin{itemize}
  \item \textbf{Claude Opus 4.5}, in Claude Code (Pro subscription, about \$20/month)
  \item \textbf{gpt-5.2-codex (high)}, in Codex CLI (Plus subscription, about \$20/month)
\end{itemize}
This guide works for both the CLI versions (terminal) and the VSCode extension versions (Codex and Claude Code).

\textit{Note: Earlier versions of this guide used Grok 3, then Gemini 2.5 Pro. It now uses Claude Opus 4.5 or gpt-5.2-codex (high).}

\textit{Note 2: If you want to use Cursor, refer to version 1.1 of this guide (\href{https://github.com/EnzeD/vibe-coding/tree/1.1.1}{link}). We believe Codex CLI or Claude Code is more powerful.}

Setting things up correctly is key. If you're serious about creating a fully functional and visually appealing game (or app), take the time to build a solid foundation.

\section{Setting Up Everything}
\subsection{0. Initial Setup}
\begin{itemize}
  \item Download Visual Studio Code: \href{https://code.visualstudio.com/}{https://code.visualstudio.com/}
  \item Create a new folder on your machine
  \item Right-click the folder and open it with VS Code
  \item Start a terminal
  \begin{itemize}
    \item For Claude Code: \texttt{curl -fsSL https://claude.ai/install.sh | bash}
    \item For Codex CLI: install Node, then \texttt{npm i -g @openai/codex}
  \end{itemize}
\end{itemize}

\subsection{1. Game Design Document (or PRD for an app)}
\begin{itemize}
  \item Ask GPT-5.2 or Opus 4.5 to create a simple \textbf{Game Design Document} in Markdown: \texttt{game-design-document.md}.
  \item Review and refine the document to ensure it matches your vision. Keep it simple; the goal is to give your AI context about structure and intent.
  \item If you want, have the AI ask you questions and use that to write the GDD or PRD.
\end{itemize}

\subsection{2. Tech Stack and \texttt{CLAUDE.md} / \texttt{Agents.md}}
\begin{itemize}
  \item Ask GPT-5.2 or Opus 4.5 for the best stack for your game (for example, Vite + ThreeJS + WebSocket for a multiplayer 3D game). Save it as \texttt{tech-stack.md}.
  \item Challenge it to propose the simplest yet most robust stack possible.
  \item In your terminal, open Claude Code or Codex CLI and use the \texttt{/init} command. It will use the two Markdown files you created. This generates rules that guide your LLM.
  \item Crucially, review the generated rules. Make sure they emphasize modularity (multiple files) and discourage a monolith. Adjust rules and triggers if needed.
  \item Some rules should be marked as \textbf{``Always''}. Examples:
\end{itemize}
\begin{verbatim}
# IMPORTANT:
# Always read memory-bank/@architecture.md before writing any code. Include entire database schema.
# Always read memory-bank/@game-design-document.md before writing any code.
# After adding a major feature or completing a milestone, update memory-bank/@architecture.md.
\end{verbatim}
\begin{itemize}
  \item Ensure other rules guide best practices for your stack (networking, state management, etc.). This setup is mandatory for clean, optimized code.
\end{itemize}

\subsection{3. Implementation Plan}
\begin{itemize}
  \item Provide GPT-5.2 or Opus 4.5 with \texttt{game-design-document.md} and \texttt{tech-stack.md}.
  \item Ask for a detailed \textbf{Implementation Plan} in Markdown: step-by-step instructions for your AI developers.
  \item Steps should be small and specific, and each step must include a test to validate correctness.
  \item No code in this document—only clear, concrete instructions focused on the base game (features come later).
\end{itemize}

\subsection{4. Memory Bank}
\begin{itemize}
  \item Create a project folder and open it in VS Code.
  \item Inside it, create a subfolder \texttt{memory-bank}.
  \item Add the following files:
  \begin{itemize}
    \item \texttt{game-design-document.md}
    \item \texttt{tech-stack.md}
    \item \texttt{implementation-plan.md}
    \item \texttt{progress.md} (empty; for tracking completed steps)
    \item \texttt{architecture.md} (empty; for documenting file purposes)
  \end{itemize}
\end{itemize}

\section{Vibe Coding the Base Game}
\subsection{Make Sure Everything Is Clear}
\begin{itemize}
  \item In VS Code, open Codex or Claude Code (extension or terminal).
  \item Prompt: Read all documents in \texttt{/memory-bank}. Is \texttt{implementation-plan.md} clear? What questions do you have to make it 100\% clear?
  \item The AI usually asks 9--10 questions. Answer them and have it edit \texttt{implementation-plan.md} accordingly.
\end{itemize}

\subsection{Your First Implementation Prompt}
\begin{itemize}
  \item Open Codex or Claude Code in VS Code or the terminal.
  \item Prompt: Read all documents in \texttt{/memory-bank}, and proceed with Step 1 of the implementation plan. I will run the tests. Do not start Step 2 until I validate the tests. Once I validate them, open \texttt{progress.md} and document what you did for future developers. Then add any architectural insights to \texttt{architecture.md} to explain what each file does.
  \item Always start with Ask mode or Plan Mode, and once you are satisfied, allow the AI to proceed.
  \item Extreme vibe: Install \href{https://superwhisper.com}{Superwhisper} to speak casually with Claude or GPT-5.2 instead of typing.
\end{itemize}

\subsection{Workflow}
\begin{itemize}
  \item After completing Step 1, commit your changes to Git.
  \item Start a new chat (\texttt{/new} or \texttt{/clear}). Large context windows improve output quality.
  \item Prompt: Go through all files in \texttt{memory-bank}, read \texttt{progress.md} to understand prior work, and proceed with Step 2. Do not start Step 3 until I validate the test.
  \item Repeat until the entire \texttt{implementation-plan.md} is complete.
\end{itemize}

\section{Adding Details}
After the base game is built, add features iteratively.
\begin{itemize}
  \item Want fog, post-processing, effects, or sounds? Better models or art assets? A sky upgrade?
  \item For each major feature, create a new \texttt{feature-implementation.md} with short steps and tests.
  \item Implement and test incrementally.
\end{itemize}

\section{Fixing Bugs and Stuckness}
\begin{itemize}
  \item If a prompt fails or breaks the game, use \texttt{/rewind} in Claude Code and refine your prompt. In GPT-5.2, commit often and reset as needed.
  \item For errors:
  \begin{itemize}
    \item If JavaScript: open the browser console (F12), copy the error, and paste it into VS Code. Provide a screenshot for visual glitches.
    \item Lazy option: install \href{https://browsertools.agentdesk.ai/installation}{BrowserTools} to avoid manual copying or screenshots.
  \end{itemize}
  \item If stuck: revert to your last Git commit and retry with new prompts.
  \item If really stuck: use \href{https://repoprompt.com/}{RepoPrompt} or \href{https://uithub.com/}{uithub} to view your whole codebase and ask GPT-5.2 or Claude for help.
\end{itemize}

\section{Claude Code and Codex Tips}
\begin{itemize}
  \item Run Codex CLI or Claude Code in the VS Code terminal to view diffs and feed context without leaving your workspace.
  \item Claude Code \texttt{/rewind} lets you roll the project back if an iteration misses the mark.
  \item Custom commands or skills can speed you up (for example, \texttt{/explain \$arguments} to trigger a deep dive).
  \item Clear context frequently with \texttt{/clear} or \texttt{/compact}; keep context usage under 50--60\% for best performance.
  \item Save time at your own risk: \texttt{claude --dangerously-skip-permissions} or \texttt{codex --yolo} starts in a no-confirmation mode.
\end{itemize}

\section{Other Tips}
\begin{itemize}
  \item Small edits: use GPT-5.2 (medium).
  \item Great marketing copy: use Opus 4.5.
  \item Generate sprites (2D images): use ChatGPT and Nano Banana Pro.
  \item Generate 3D assets: use Trellis, Tripo, or Hunyuan.
  \item Generate music: use Suno or ElevenLabs.
  \item Generate sound effects: use ElevenLabs.
  \item Generate video: use Sora 2 or Veo 3.
  \item Better prompt outputs: add ``think as long as needed to get this right; I am not in a hurry. What matters is that you follow precisely what I ask you and execute it perfectly. Ask me questions if I am not precise enough.'' For Claude Code, deeper reasoning cues: \texttt{think} $<$ \texttt{think hard} $<$ \texttt{think harder} $<$ \texttt{ultrathink}.
\end{itemize}

\section{Frequently Asked Questions}
\begin{description}[leftmargin=0pt]
  \item[Projects vibe coded 100\% with this method?] Recent examples: \url{https://fly.zullo.fun/} (3D WW2 dogfight arena), \url{https://vibecraft.game/} (prompt-to-existence 3D game), \url{https://www.dow-de.com/} (Warhammer 40000 Dawn of War stats web app).
  \item[I'm building an app, not a game. Same workflow?] Mostly yes. Use a PRD instead of a GDD. Tools like v0, Lovable, or Bolt.new are fine for prototyping; move code to GitHub, then clone to VS Code or the terminal with this guide.
  \item[I can't replicate the plane in your dogfight game!] It took about 30 prompts guided by a specific \texttt{plane-implementation.md}. Use sharp prompts like ``cut out space in the wings for ailerons,'' not vague ones like ``make a plane.''
  \item[Why are Claude Code or Codex CLI better than Cursor right now?] They pair better with Claude Opus 4.5 and GPT-5.2 than Cursor does, work well in the terminal from any IDE, and support powerful customization (custom commands, sub-agents, hooks). Lower-tier Claude or ChatGPT plans are enough to start.
  \item[I don't know how to set up a server for my multiplayer game.] Ask your AI.
\end{description}

\end{document}
